% ============================================================
% insta_sialkot_football_marx_SQUARECANVAS_BLACK_BOXED.tex
% Standalone (pdflatex). Each page = one square slide.
% Black background, white text, white box outline.
% Asymmetric margins (more space to right and bottom).
% Header font size controlled independently from body font.
% PKR first everywhere; USD in brackets using a stated reference rate.
% ============================================================

\documentclass[17pt]{extarticle}

% --- Square canvas; more usable area toward right + bottom ---
\usepackage[
  paperwidth=7.5in,paperheight=7.5in,
  left=0.33in,top=0.33in,
  right=0.22in,bottom=0.22in
]{geometry}

\usepackage[T1]{fontenc}
\usepackage{lmodern}
\usepackage{microtype}

\usepackage{xcolor}
\usepackage{amsmath,amssymb}
\usepackage{siunitx}
\usepackage{enumitem}
\usepackage[hidelinks]{hyperref}
\usepackage{xfp}

\usepackage{xurl} % better URL line breaking than plain \url
\urlstyle{same}

\sisetup{detect-all,group-separator={,},group-minimum-digits=4}

% --- Black canvas, white ink ---
\AtBeginDocument{\pagecolor{black}\color{white}}

% --- Spacing (looser, readable on phone) ---
\setlength{\parindent}{0pt}
\setlength{\parskip}{0.60em}
\renewcommand{\baselinestretch}{1.12}

\setlist[itemize]{leftmargin=1.35em,itemsep=0.28em,topsep=0.20em,parsep=0pt}

% --- Robust line breaking inside boxed slides ---
\emergencystretch=2.5em

% --- Units (avoid \day/\month primitives) ---
\newcommand{\uDay}{\mathrm{day}}
\newcommand{\uHour}{\mathrm{hour}}
\newcommand{\uBall}{\mathrm{ball}}
\newcommand{\uMonth}{\mathrm{month}}

% --- Currency tokens and formats ---
\newcommand{\PKR}{\text{PKR}}
\newcommand{\USD}{\text{USD}}
\newcommand{\pkr}[1]{\text{PKR}\,\num{#1}}
\newcommand{\usd}[1]{\text{USD}\,\num{#1}}

% ============================================================
% Exchange rate (single point of truth)
% ============================================================
\newcommand{\USDPKR}{280.10}
\newcommand{\RateDate}{2 Jan 2026}

\newcommand{\USDofPKR}[1]{\fpeval{round((#1)/\USDPKR,2)}}
\newcommand{\PKRofUSD}[1]{\fpeval{round((#1)*\USDPKR,0)}}
\newcommand{\PKRofUSDtwo}[1]{\fpeval{round((#1)*\USDPKR,2)}}

% ============================================================
% Fixed inputs (your numbers; PKR-side is “as given”)
% ============================================================

% Piece-rate example:
\newcommand{\PieceRatePKR}{160}
\newcommand{\BallsPerDay}{3}
\newcommand{\HoursPerBall}{3}

\newcommand{\DailyPayPKR}{\fpeval{round(\PieceRatePKR*\BallsPerDay,0)}}   % 480
\newcommand{\DailyHours}{\fpeval{round(\HoursPerBall*\BallsPerDay,0)}}    % 9
\newcommand{\HourlyWagePKR}{\fpeval{round(\DailyPayPKR/\DailyHours,2)}}   % 53.33

\newcommand{\MonthlyPayPKR}{9600}
\newcommand{\LivingWagePKR}{20000}
\newcommand{\ShortfallPKR}{\fpeval{round(\LivingWagePKR-\MonthlyPayPKR,0)}}
\newcommand{\LivingPct}{\fpeval{round(100*\MonthlyPayPKR/\LivingWagePKR,0)}}

% Producer cost structure example (PKR per ball):
\newcommand{\cPKR}{154.52}
\newcommand{\vPKR}{56.29}
\newcommand{\kPKR}{210.83}
\newcommand{\MarkupPct}{8.1}
\newcommand{\pPKR}{227.91}
\newcommand{\sPKR}{17.08}
\newcommand{\NewValuePKR}{73.37}

% International benchmarks (USD -> PKR via reference rate):
\newcommand{\RetailUSD}{165}
\newcommand{\RetailPKR}{\PKRofUSD{\RetailUSD}}

\newcommand{\ExportAvgUSD}{4.44}
\newcommand{\ExportAvgPKR}{\PKRofUSDtwo{\ExportAvgUSD}}
\newcommand{\ExportAvgPKRint}{\fpeval{round(\ExportAvgPKR,0)}}

% Derived comparisons:
\newcommand{\DailyExportPKR}{\fpeval{round(\BallsPerDay*\ExportAvgPKR,0)}}
\newcommand{\DailyRetailPKR}{\fpeval{round(\BallsPerDay*\RetailPKR,0)}}
\newcommand{\WageShareExportPct}{\fpeval{round(100*\DailyPayPKR/(\BallsPerDay*\ExportAvgPKR),2)}}
\newcommand{\WageShareRetailPct}{\fpeval{round(100*\DailyPayPKR/(\BallsPerDay*\RetailPKR),3)}}

\newcommand{\HoursToBuy}{\fpeval{round(\RetailPKR/\HourlyWagePKR,0)}}
\newcommand{\BallsToBuy}{\fpeval{round((\RetailPKR/\HourlyWagePKR)/\HoursPerBall,0)}}
\newcommand{\DaysToBuy}{\fpeval{round((\RetailPKR/\HourlyWagePKR)/\DailyHours,1)}}

% Time intuition:
\newcommand{\NecessaryHours}{6.9}
\newcommand{\SurplusHours}{2.1}

% ============================================================
% Typography controls (HEADER vs BODY)
% ============================================================
\newcommand{\BodyFont}{\normalsize} % follows document base (17pt)
\newcommand{\HeaderFont}{\fontsize{20}{24}\selectfont}
\newcommand{\SlideNoFont}{\fontsize{16}{19}\selectfont}

% ============================================================
% Boxed slide macro (white outline; vertically centered body)
% ============================================================

\setlength{\fboxsep}{14pt}
\setlength{\fboxrule}{1.05pt}

\newlength{\InnerW}
\newlength{\InnerH}
\setlength{\InnerW}{\dimexpr\textwidth-2\fboxsep-2\fboxrule\relax}
\setlength{\InnerH}{\dimexpr\textheight-2\fboxsep-2\fboxrule\relax}

\newcommand{\HeaderLine}[2]{%
  \parbox[t]{0.78\linewidth}{\raggedright{\HeaderFont\bfseries #1}}%
  \hfill
  \parbox[t]{0.20\linewidth}{\raggedleft{\SlideNoFont\bfseries #2}}%
  \par\vspace{0.25em}\hrule\vspace{0.60em}%
}

\newcommand{\SourceLine}[1]{%
  \vspace{0.55em}
  {\footnotesize\raggedright\textit{Source/benchmark:} #1\par}
}

\newcommand{\SlideBox}[3]{%
  \clearpage
  \thispagestyle{empty}
  \noindent\fbox{%
    \begin{minipage}[t][\InnerH][t]{\InnerW}
      \HeaderLine{#1}{#2}
      \vspace*{\fill}
      {\BodyFont\raggedright\sloppy #3\par}
      \vspace*{\fill}
    \end{minipage}
  }%
}

\begin{document}

% ============================================================
% SLIDES
% ============================================================

\SlideBox
{Sialkot footballs: workshop labour, world-market prices}
{01/12}
{%
\textbf{Reference rate used here:} \(1~\USD \approx \pkr{\USDPKR}\) (mid-market, \RateDate). We write PKR first and show USD in brackets.

Sialkot’s football workers produce commodities that pass through exporters, brands, and retailers into global leagues and tournaments. The same stitched object can reappear in the world market at a money price that stands far above the local wage.

Two benchmarks (converted using the reference rate above):
\[
\begin{array}{ll}
\text{Retail benchmark:} & \pkr{\RetailPKR}\ (\usd{\RetailUSD})\\
\text{Export-average benchmark:} & \pkr{\ExportAvgPKRint}\ (\usd{\ExportAvgUSD})
\end{array}
\]

We will measure one simple thing: how wages, costs, and surplus fit together, and what that implies for class struggle in a supply chain.%
}

% --- ADD THIS SLIDE AFTER CURRENT SLIDE 01 ---

\SlideBox
{Data sources:}
{02/12}
{%
\small % <-- ADD THIS LINE (first line inside the body)

These slides combine figures drawn from different sources and years (each labelled). That is deliberate: the aim is to make the class relation legible, not to pretend all numbers are a single same-year accounts sheet.

\begin{itemize}
\item \textbf{Piece-rate, time per ball, and the 9,600/month example:} reported in \textbf{Nov 2022} reporting on Sialkot stitchers (piece-rate \pkr{160}, \(\sim 3\) hours/ball, \(\sim 3\) balls/day).
\item \textbf{“Living wage \(\sim\) \pkr{20000}/month”:} quoted as a \textbf{researcher-estimate benchmark} in the same \textbf{Nov 2022} reporting (we update this benchmark below).
\item \textbf{Export scale:} \textbf{FY 2021/22} total output/value reported by the local chamber of commerce (via Reuters): 43 million balls valued at USD 191 million (implied average export unit value \(\approx \usd{4.44}\)).
\item \textbf{Producer-stage cost structure \(c,v,k\) and profit margin:} baseline factory survey for Sialkot footballs in \textbf{2012} (published as an NBER working paper in \textbf{2015}).
\item \textbf{Living wage update used here:} \textbf{2025 estimate} (Anker Method; GLWC/Anker Research Institute).
\item \textbf{FX conversion shown in brackets:} SBP USD/PKR M2M revaluation rate as on \RateDate.
\end{itemize}

We keep PKR as primary, and we mark benchmark-years so inflation and time are not hidden.%
}


\SlideBox
{Notations first:}
{03/12}
{%
We use standard Marxian notation, defined before any calculation:

\begin{itemize}
\item \(w_{\uBall}\): piece-rate pay per ball. \(w_{\uDay}\), \(w_{\uMonth}\): wage income per day/month.
\item \(T_{\uDay}\): hours worked per day. \(w_{\uHour}=w_{\uDay}/T_{\uDay}\): implied hourly wage.
\item \(c\): \textbf{constant capital} (materials, overhead, machines) whose value is transferred.
\item \(v\): \textbf{variable capital} (wages paid for labour-power).
\item \(k=c+v\): \textbf{cost price} advanced by capital.
\item \(p\): selling price at the producer stage. \(s=p-k\): \textbf{surplus} in money form at this stage.
\item \(e=\dfrac{s}{v}\): \textbf{rate of exploitation} (surplus relative to wages).
\end{itemize}

From here on, each slide repeats the definition relevant to its step, and then plugs in numbers.%
}

\SlideBox
{Piece rate \(\rightarrow\) hourly wage: the wage in disguise}
{04/12}
{%
\textbf{Definitions.} \(w_{\uBall}\) is pay per ball; \(q\) is balls/day; \(t_{\uBall}\) is hours/ball.
Then \(w_{\uDay}=w_{\uBall}q\), \(T_{\uDay}=t_{\uBall}q\), and \(w_{\uHour}=w_{\uDay}/T_{\uDay}\).

Numbers (PKR first; USD in brackets):
\[
w_{\uBall}=\pkr{\PieceRatePKR}\ (\usd{\USDofPKR{\PieceRatePKR}}),\quad
q=\BallsPerDay,\quad
t_{\uBall}=\HoursPerBall.
\]
Daily wage income:
\[
w_{\uDay}=160\times 3=\pkr{\DailyPayPKR}\ (\usd{\USDofPKR{\DailyPayPKR}}).
\]
Daily hours:
\[
T_{\uDay}=3\times 3=\DailyHours\ \text{hours}.
\]
Implied hourly wage:
\[
w_{\uHour}=\frac{480}{9}\approx \pkr{\HourlyWagePKR}/\uHour\ (\usd{\USDofPKR{\HourlyWagePKR}}/\uHour).
\]

A piece rate looks like payment for output. In reality it is a method of controlling time: the worker must increase pace, lengthen hours, or absorb risk to keep income steady.%
}

\SlideBox
{Monthly wage vs. survival: social reproduction in numbers}
{05/12}
{%
\textbf{Definitions.} \(w_{\uMonth}\) is monthly wage income. \(L_{\uMonth}\) is a living-cost benchmark (the monthly cost to reproduce labour-power).

Reported monthly income at this pace:
\[
w_{\uMonth}=\pkr{\MonthlyPayPKR}\ (\usd{\USDofPKR{\MonthlyPayPKR}}).
\]
Living-cost benchmark:
\[
L_{\uMonth}=\pkr{\LivingWagePKR}\ (\usd{\USDofPKR{\LivingWagePKR}}).
\]
Coverage and shortfall:
\begin{align*}
\frac{w_{\uMonth}}{L_{\uMonth}}
&=\frac{9600}{20000}=0.48 \Rightarrow \LivingPct\%,\\[0.25em]
L_{\uMonth}-w_{\uMonth}
&=\pkr{\ShortfallPKR}\;(\usd{\USDofPKR{\ShortfallPKR}}).
\end{align*}


This is social reproduction as arithmetic: when wages sit below reproduction, the “adjustment” is rationing, debt, longer hours, and unpaid labour pushed into households.%
}

\SlideBox
{What the worker produces in a day}
{06/12}
{%
\textbf{Definitions.} Output per day is \(q\). Daily wage is \(w_{\uDay}\). We compare the wage to two money benchmarks for the day’s output.

At this pace \(q=3\) balls/day and \(w_{\uDay}=\pkr{\DailyPayPKR}\ (\usd{\USDofPKR{\DailyPayPKR}})\).

Benchmark values of the day’s output:
\begin{align*}
\text{Export benchmark/day}
&=3\times \pkr{\ExportAvgPKRint}\approx \pkr{\DailyExportPKR},\\
\text{Retail benchmark/day}
&=3\times \pkr{\RetailPKR}= \pkr{\DailyRetailPKR}.
\end{align*}

Wage as a share of those benchmarks:
\begin{align*}
\frac{w_{\uDay}}{\text{export benchmark/day}}
&=\frac{480}{\num{\DailyExportPKR}}\approx \WageShareExportPct\%,\\
\frac{w_{\uDay}}{\text{retail benchmark/day}}
&=\frac{480}{\num{\DailyRetailPKR}}\approx \WageShareRetailPct\%.
\end{align*}

These ratios do not claim the local employer pockets full retail. They show the scale of separation between the wage and the commodity’s world-market money form.%
}

\SlideBox
{Inside one ball: \(c\), \(v\), \(k\)}
{07/12}
{%
\textbf{Definitions.}
\(c\) = constant capital \emph{advanced and used up} in production (value transferred);
\(v\) = variable capital (wages paid for labour-power);
\(k=c+v\) = cost price advanced by capital.

\textbf{Breakdown of \(c\).} In general,
\[
c = c_{\text{materials}} + c_{\text{depreciation}} + c_{\text{overhead}},
\]
where materials include leather/synthetic panels, bladder, thread, glue, packaging; depreciation is wear-and-tear of tools/machines/buildings; overhead includes electricity, rent, admin, and related running costs. In the producer-cost figures used here, these elements are reported together as one combined \(c\) (“materials + overhead”), so we do not split them numerically.

Representative producer cost structure (PKR per ball):
\[
c=\pkr{\cPKR},\qquad
v=\pkr{\vPKR},\qquad
k=c+v=\pkr{\kPKR}.
\]


Why it matters: \(c\) does not create new value; it passes value on. New value comes from living labour. Profit is not explained by “exchange” alone; it is rooted in production under capitalist control.%
}

\SlideBox
{Selling price \(p\), surplus \(s\), and new value \(v+s\)}
{08/12}
{%
\textbf{Definitions.}
\(p\) = producer-stage selling price;
\(s=p-k\) = surplus in money form at this stage;
new value (money form) \(=v+s\).

With markup \(m=\MarkupPct\%\) on cost:
\[
p=k\Bigl(1+\frac{m}{100}\Bigr)\approx \pkr{\pPKR}.
\]
So surplus:
\[
s\approx p-k=\pkr{\pPKR}-\pkr{\kPKR}=\pkr{\sPKR}.
\]
New value added by living labour (money form):
\[
v+s=\pkr{\vPKR}+\pkr{\sPKR}=\pkr{\NewValuePKR}.
\]

Labour creates new value; wages return only part of it. The remainder appears as \(s\), and that remainder is what capital lives on.%
}

\SlideBox
{Rate of exploitation \(e=s/v\)}
{09/12}
{%
\textbf{Definition.} The rate of exploitation is \(e=\dfrac{s}{v}\): surplus relative to wages.

Compute:
\[
e=\frac{17.08}{56.29}\approx 0.303 \Rightarrow 30.3\%.
\]
Meaning: for every \(\pkr{100}\) paid in wages, about \(\pkr{30}\) is appropriated as surplus at this stage.

Time intuition (for a 9-hour day):
\[
T_{\uDay}=T_n+T_s \approx \NecessaryHours + \SurplusHours = 9,
\]
where \(T_n\) is necessary labour time (reproduces the wage) and \(T_s\) is surplus labour time (appropriated by capital).%
}

\SlideBox
{How long to buy what you make?}
{10/12}
{%
\textbf{Definitions.} \(w_{\uHour}\) is hourly wage. \(H=\dfrac{\text{price}}{w_{\uHour}}\) is hours needed to earn a given price in wages.

Hourly wage:
\[
w_{\uHour}\approx \pkr{\HourlyWagePKR}/\uHour\ (\usd{\USDofPKR{\HourlyWagePKR}}/\uHour).
\]
Retail benchmark price:
\[
\pkr{\RetailPKR}\ (\usd{\RetailUSD}).
\]
Hours to afford one retail-price ball:
\[
H=\frac{\num{\RetailPKR}}{\num{\HourlyWagePKR}}\approx \HoursToBuy\ \text{hours}.
\]
Convert into output and workdays:
\[
\HoursToBuy\div 3 \approx \BallsToBuy\ \text{balls},
\qquad
\HoursToBuy\div 9 \approx \DaysToBuy\ \text{days}.
\]

So the worker must stitch roughly \textbf{\BallsToBuy\ balls}—nearly \textbf{\DaysToBuy\ workdays}—to earn enough wages to buy \textbf{one} ball at the retail benchmark.%
}

\SlideBox
{Why Marxism clarifies what liberal economics obscures}
{11/12}
{%
Neoclassical and neoliberal economics typically start from exchange: a wage contract is made, a price is paid, and outcomes are explained as “returns” to factors.

Marxism starts from production and power. Workers sell \emph{labour-power} for wages \(v\). Under capitalist command, labour produces new value \(v+s\). Because capital owns the product, it appropriates \(s\). Exploitation is the measurable split:
\[
\text{new value}=v+s,\qquad e=\frac{s}{v}.
\]

This is why exploitation can exist even when contracts are legal and wages are paid on time. The issue is not merely a defective contract; it is the property relation that makes the contract necessary.%
}

\SlideBox
{History as class struggle:}
{12/12}
{%
When Marx and Engels say history is the history of class struggles, they mean the central conflict is over control of labour, ownership of the product, and appropriation of surplus.

In a supply chain like football manufacturing, that conflict appears as:
\begin{itemize}
\item piece rates and speed-up (how time is commanded for how little pay),
\item subcontracting and fragmentation (how capital isolates workers),
\item compliance regimes and buyer power (how global contracts discipline local labour),
\item reproduction crises (how wages below survival are “managed” through debt and unpaid labour).
\end{itemize}

The arithmetic is not decoration. It makes the relation legible: who bears the cost, who captures the surplus, and why collective organisation becomes necessary.%
}

\SlideBox
{Perspectives for the movement}
{13}
{%
If wages are below reproduction (\pkr{\MonthlyPayPKR} vs \pkr{\LivingWagePKR}), closing the reproduction gap is not charity; it is a class demand. If surplus exists and is measurable (\(e\approx 30.3\%\) in the producer-stage example), the struggle is also over time, intensity, and control.

Practical implications:
\begin{itemize}
\item Raise piece rates and base wages so monthly income meets and exceeds the reproduction benchmark.
\item Reduce working time without wage loss: contest surplus labour time directly.
\item Organise across workshops and subcontractors so capital cannot isolate workers by contract type.
\item Target downstream enforcement: purchasing contracts and brand/retail power that squeezes labour at the bottom.
\end{itemize}

Marxist rigour is useful because it ties demands to the structure of the relation: \(v\) versus \(v+s\), reproduction versus deprivation, and the necessity of collective power to change the terms.%
}

\end{document}
